\chapter{Results}
\label{ch:results}

This chapter reports the final confirmatory evaluation described in Chapter~\ref{chap:methods}. The analysis was restricted to the matched 1437-complex subset for which both the full and interface graph views were available after molecular-dynamics processing, force-feature construction, and interface materialization. The final model matrix comprised 52 configurations: 48 GraphVAE-based representation models and 4 direct supervised baselines. Each configuration was evaluated across 10 outer held-out fits obtained from 5 folds and 2 repeats. If $m_{c,t}$ denotes a metric for configuration $c$ on outer fit $t \in \{1,\dots,10\}$, the reported configuration-level summary is
\begin{equation}
\bar{m}_{c} = \frac{1}{10}\sum_{t=1}^{10} m_{c,t}.
\end{equation}
All rankings and factor comparisons below are based on these configuration-level means. For paired factor comparisons, descriptive 95\% intervals over paired configuration deltas are reported as $\bar{\Delta} \pm 1.96\, s_{\Delta}/\sqrt{n}$, where $n$ is the number of paired configurations. These intervals are included to show the stability of the direction and magnitude of the observed effect; they are not formal hypothesis tests.

\section{Headline model performance}

Table~\ref{tab:headline-models} lists the strongest configurations from the final matrix. The best overall model was the direct supervised baseline on the full static graph, \texttt{full\_base\_S}, with mean test RMSE 2.3379, mean test Pearson correlation 0.5064, and mean test $R^2$ 0.2146. The best GraphVAE configuration was \texttt{full\_S\_semi\_shared\_static\_z32}, with mean test RMSE 2.4292, mean test Pearson correlation 0.4649, and mean test $R^2$ 0.1537. The best static-plus-dynamic GraphVAE configuration was \texttt{full\_SD\_semi\_shared\_static\_z8}, with mean test RMSE 2.4364 and mean test Pearson correlation 0.4749. The best interface-only models were materially weaker than their full-graph counterparts. The strongest interface baseline, \texttt{iface\_base\_S}, achieved mean test RMSE 2.5782, while the best interface GraphVAE, \texttt{iface\_S\_semi\_shared\_static\_z8}, achieved mean test RMSE 2.6041.

\begin{table}[t]
\centering
\small
\begin{tabular}{|p{6.2cm}|c|c|c|c|}
\hline
Configuration & Test RMSE & Test $r$ & Test $R^2$ & Test$-$Train RMSE gap \\
\hline
\texttt{full\_base\_S} & 2.3379 & 0.5064 & 0.2146 & 0.8021 \\
\texttt{full\_base\_SD} & 2.4073 & 0.4802 & 0.1685 & 0.8414 \\
\texttt{full\_S\_semi\_shared\_static\_z32} & 2.4292 & 0.4649 & 0.1537 & 0.7451 \\
\texttt{full\_SD\_semi\_shared\_static\_z8} & 2.4364 & 0.4749 & 0.1478 & 1.1303 \\
\texttt{iface\_base\_S} & 2.5782 & 0.3451 & 0.0470 & 0.8090 \\
\texttt{iface\_S\_semi\_shared\_static\_z8} & 2.6041 & 0.3402 & 0.0282 & 0.8018 \\
\texttt{iface\_SD\_semi\_sd\_static\_plus\_dynamic\_all\_z16} & 2.6238 & 0.3252 & 0.0146 & 0.7597 \\
\hline
\end{tabular}
\caption{Selected headline models from the final 52-configuration confirmatory matrix. Metrics are means over 10 outer held-out fits on the matched 1437-complex set. For GraphVAE models, the reported predictive metrics come from the downstream ridge probe on latent means rather than solely from the neural affinity head.}
\label{tab:headline-models}
\end{table}

The best direct supervised baseline outperformed the best GraphVAE by 0.0913 RMSE units and by 0.0415 Pearson-correlation units on average. Within the $SD$ family, the best model was again a full-graph method rather than an interface method: \texttt{full\_base\_SD} reached mean test RMSE 2.4073 and mean test $r=0.4802$, while the best $SD$ GraphVAE reached mean test RMSE 2.4364 and mean test $r=0.4749$. The confirmatory matrix therefore did not support a claim that the variational representation-learning objective outperformed a direct supervised graph regressor on the present dataset and feature set.

\section{Effect of graph view: full versus interface}

The clearest factor effect in the matrix was the advantage of the full graph view over the interface graph view. There were 24 matched GraphVAE configuration pairs differing only in whether the model was trained on \texttt{full} or \texttt{interface} records. In all 24 of these paired comparisons, the full graph achieved a lower mean test RMSE than the corresponding interface graph. Defining the paired RMSE penalty for the interface representation as
\begin{equation}
\Delta_{\mathrm{RMSE}}^{\mathrm{iface-full}} = \bar{\mathrm{RMSE}}_{\mathrm{iface}} - \bar{\mathrm{RMSE}}_{\mathrm{full}},
\end{equation}
the mean paired penalty was 0.1005 RMSE units, with a descriptive 95\% interval of [0.0653, 0.1356]. If the corresponding Pearson-correlation advantage is defined as
\begin{equation}
\Delta_{r}^{\mathrm{full-iface}} = \bar{r}_{\mathrm{full}} - \bar{r}_{\mathrm{iface}},
\end{equation}
the mean full-graph advantage was 0.0993, with a descriptive 95\% interval of [0.0847, 0.1138].

The same direction held at the family level. Averaging over all 24 full-graph GraphVAE configurations gave mean test RMSE 2.5494 and mean test Pearson correlation 0.2883. Averaging over all 24 interface GraphVAE configurations gave mean test RMSE 2.6498 and mean test Pearson correlation 0.1890. The direct supervised baselines showed the same pattern. The best full baseline, \texttt{full\_base\_S}, achieved mean test RMSE 2.3379, whereas the best interface baseline, \texttt{iface\_base\_S}, achieved mean test RMSE 2.5782. This confirms that the superiority of the full representation was not specific to the variational model family.

These results are sufficiently uniform that the main view-level conclusion is not sensitive to model-selection details. On the present matched subset, with the present residue graph definitions and interface-patch construction, the full-graph representation was consistently more predictive than the interface-only representation.

\section{Effect of dynamic features: $S$ versus $SD$}

The second main question was whether augmenting the static graph with molecular-dynamics-derived features improved predictive performance. The cleanest comparison for this purpose is between matched $S$ and $SD$ configurations under the \texttt{shared\_static} reconstruction target, because both models are then trained under the same target policy and differ only in whether the dynamic channels are included as inputs. There were 12 such matched GraphVAE pairs in the final matrix. In 11 of the 12 pairs, the $S$ model achieved a lower mean test RMSE than the corresponding $SD$ model.

Defining the paired dynamic-feature penalty as
\begin{equation}
\Delta_{\mathrm{RMSE}}^{\mathrm{SD-S}} = \bar{\mathrm{RMSE}}_{\mathrm{SD}} - \bar{\mathrm{RMSE}}_{\mathrm{S}},
\end{equation}
the mean penalty was 0.0254 RMSE units, with a descriptive 95\% interval of [0.0134, 0.0375]. The corresponding Pearson-correlation difference,
\begin{equation}
\Delta_{r}^{\mathrm{S-SD}} = \bar{r}_{\mathrm{S}} - \bar{r}_{\mathrm{SD}},
\end{equation}
favored the static-only representation by 0.0215 on average, with a descriptive 95\% interval of [0.0043, 0.0387]. The direct supervised baselines were consistent with the same conclusion: \texttt{full\_base\_SD} was worse than \texttt{full\_base\_S} by 0.0694 RMSE units, and \texttt{iface\_base\_SD} was not better than \texttt{iface\_base\_S}.

This effect was small rather than catastrophic. The best $SD$ GraphVAE, \texttt{full\_SD\_semi\_shared\_static\_z8}, was only 0.0072 RMSE units worse than the best $S$ GraphVAE, \texttt{full\_S\_semi\_shared\_static\_z32}, and it had a slightly higher test Pearson correlation (0.4749 versus 0.4649). The problem is therefore not that the dynamic features always destroy performance. The problem is that, across the full confirmatory matrix, the dynamic augmentation does not yield a robust and repeatable generalization gain over the static representation.

\section{Effect of supervision}

The influence of semi-supervision depended strongly on graph view. Within the full graph view, semi-supervision helped consistently. For the static mode, all 3 of the 3 matched full-graph pairs favored semi-supervision over the corresponding unsupervised model, with a mean test-RMSE gain of 0.1640. For the static-plus-dynamic mode, all 9 of the 9 matched full-graph pairs favored semi-supervision, with a mean test-RMSE gain of 0.1604. The strongest full-graph models therefore all came from the semi-supervised branch of the matrix.

The interface results were weaker and less coherent. In the static mode, semi-supervision still helped in all 3 of the 3 matched interface pairs, but the mean gain was only 0.0290 RMSE units. In the dynamic mode, semi-supervision did not produce a robust benefit. Across the 9 matched interface $SD$ pairs, the mean semi-supervised gain was $-0.0081$ RMSE units and only 3 of the 9 pairs favored the semi-supervised model. The interpretation is therefore not simply that supervision helps universally. Rather, supervision helps when the graph representation contains sufficient predictive structure, which appears to be true for the full graphs and only marginally true for the interface graphs.

\section{Effect of reconstruction target policy}

Within the $SD$ GraphVAE branch, the confirmatory matrix also tested whether the decoder should reconstruct only static targets, only dynamic targets, or both. The clearest regime in which to examine this question is the full-graph semi-supervised $SD$ family, because this was the only dynamic branch that approached competitive held-out performance. Averaging over the three latent dimensionalities gave mean test RMSE 2.4434 for \texttt{shared\_static}, 2.4795 for \texttt{sd\_static\_plus\_dynamic\_all}, and 2.5150 for \texttt{sd\_dynamic\_all}. The corresponding mean test Pearson correlations were 0.4682, 0.4335, and 0.4221, respectively.

Pairwise comparisons by latent dimensionality showed that reconstructing only dynamic targets was consistently worse than reconstructing the shared static targets. Across the three matched latent-dimensionality pairs, the mean RMSE penalty of \texttt{sd\_dynamic\_all} relative to \texttt{shared\_static} was 0.0716, with a descriptive 95\% interval of [0.0305, 0.1127]. The mean Pearson-correlation loss was 0.0461, with a descriptive 95\% interval of [0.0158, 0.0764]. The \texttt{sd\_static\_plus\_dynamic\_all} target policy was intermediate. Its mean RMSE penalty relative to \texttt{shared\_static} was 0.0360, although this comparison was less stable across latent dimensionality because only three matched pairs were available. The corresponding mean Pearson-correlation loss was still positive at 0.0347.

These results indicate that, in the only dynamic regime that was even moderately competitive, asking the decoder to reconstruct dynamic targets did not improve held-out predictive quality. The strongest dynamic models were obtained when the dynamic channels were present as inputs but the reconstruction objective remained focused on the shared static targets.

\section{Overfitting and train--validation--test behaviour}

The confirmatory matrix also yielded a clear generalization story. To quantify overfitting, define the train--test RMSE gap and train--test correlation gap for configuration $c$ as
\begin{equation}
g_{\mathrm{RMSE}}(c) = \bar{\mathrm{RMSE}}_{\mathrm{test},c} - \bar{\mathrm{RMSE}}_{\mathrm{train},c},
\end{equation}
and
\begin{equation}
g_{r}(c) = \bar{r}_{\mathrm{train},c} - \bar{r}_{\mathrm{test},c}.
\end{equation}
Large positive values of either quantity indicate stronger train-to-test divergence.

The unsupervised GraphVAE families were characterized primarily by underfitting rather than by severe overfitting. For example, the mean train, validation, and test Pearson correlations for the full unsupervised static family were 0.2486, 0.1897, and 0.2055, respectively, with a mean train--test RMSE gap of only 0.0276. The corresponding full unsupervised dynamic family had mean train, validation, and test correlations of 0.1809, 0.0967, and 0.1074. These small train--test gaps do not indicate strong robustness; they indicate that the unsupervised models did not learn much affinity signal even on their training sets.

The overfitting story emerges in the semi-supervised families. The full semi-supervised static family had mean train, validation, and test Pearson correlations of 0.7494, 0.4666, and 0.4551, respectively, with a mean train--test RMSE gap of 0.6984. This is a moderate and credible generalization gap. The full semi-supervised dynamic family fit the training data substantially more strongly, with mean train correlation 0.8218, but it did not improve held-out performance; its mean test correlation was only 0.4413 and its mean train--test RMSE gap increased to 1.0049. The interface semi-supervised dynamic family exhibited the strongest divergence of all major families, with mean train, validation, and test correlations of 0.7685, 0.3231, and 0.3189, respectively, and a mean train--test RMSE gap of 1.0140.

Table~\ref{tab:family-overfit} summarizes these family-level patterns. The best direct supervised baseline, \texttt{full\_base\_S}, combined the highest held-out performance with a more controlled generalization gap than the semi-supervised dynamic families. Its mean train, validation, and test correlations were 0.8225, 0.5337, and 0.5064, respectively, and its mean train--test RMSE gap was 0.8021. This is still a real gap, but it is materially smaller than the dynamic semi-supervised VAE families while producing better held-out accuracy.

\begin{table}[t]
\centering
\small
\begin{tabular}{|l|c|c|c|c|c|c|}
\hline
Family & Train $r$ & Val $r$ & Test $r$ & $g_r$ & Train RMSE & Test RMSE \\
\hline
Full baseline $S$ & 0.8225 & 0.5337 & 0.5064 & 0.3161 & 1.5358 & 2.3379 \\
Full baseline $SD$ & 0.8077 & 0.5191 & 0.4802 & 0.3275 & 1.5659 & 2.4073 \\
Full semi-supervised VAE $S$ & 0.7494 & 0.4666 & 0.4551 & 0.2942 & 1.7385 & 2.4369 \\
Full semi-supervised VAE $SD$ & 0.8218 & 0.4580 & 0.4413 & 0.3805 & 1.4744 & 2.4793 \\
Interface semi-supervised VAE $S$ & 0.7083 & 0.3356 & 0.3361 & 0.3722 & 1.8467 & 2.6112 \\
Interface semi-supervised VAE $SD$ & 0.7685 & 0.3231 & 0.3189 & 0.4496 & 1.6479 & 2.6620 \\
\hline
\end{tabular}
\caption{Family-level train/validation/test behaviour. Each row is the arithmetic mean over all configurations in that family, using the configuration-level means from the 10 outer held-out fits. Here $g_r = \bar{r}_{\mathrm{train}} - \bar{r}_{\mathrm{test}}$.}
\label{tab:family-overfit}
\end{table}

The worst individual overfitting case in the semi-supervised GraphVAE matrix was \texttt{iface\_SD\_semi\_shared\_static\_z32}. This model reached mean train Pearson correlation 0.8213 but only mean test Pearson correlation 0.3185, giving a train--test correlation gap of 0.5028. Its mean train RMSE was 1.4239, while its mean test RMSE was 2.6705. This specific result reinforces the broader pattern seen at the family level: adding dynamic channels, especially on the interface representation, increased training-set fit more than it increased robust held-out predictive signal.

\section{Summary of empirical conclusions}

The confirmatory matrix supports four main empirical conclusions. First, the full graph view was consistently and substantially stronger than the interface graph view for the present problem setup. Second, the current molecular-dynamics-derived feature construction did not provide a robust generalization advantage over the static representation. Third, within the dynamic GraphVAE branch, reconstruction objectives focused on dynamic targets were not helpful and were usually harmful. Fourth, the strongest overall model was the direct supervised baseline on the full static graph, not the variational model.

These conclusions are stronger than the corresponding conclusions from earlier single-split experiments because they are based on repeated outer held-out evaluation over 52 configurations and 520 total fits. At the same time, the results remain specific to the present implementation choices: short explicit-solvent trajectories, static residue graphs with time-aggregated annotations, a 128-node interface patch, and approximate protein-only inter-chain force proxies. The negative result for the current dynamic encoding should therefore be interpreted as a result about this particular representation strategy rather than as a universal statement that protein dynamics are irrelevant for affinity prediction.
